%\documentclass[DiscStateSpaceToolkit.tex]{subfiles} 
%\begin{document}
 
Hugget (1993) models a pure exchange economy to look at the question of how idiosyncratic shocks, incomplete markets, 
and borrowing constraints interact and their possible effects on interest rates.

In this model agents agents aim to maximize their consumption by saving \& borrowing (pure-exchange) assets (ie. by lending
and borrowing from one another). The amount they can borrow is limited (to roughly one years income, $\underline{a}=-5.3$; the
model is calibrated to six-periods per year). In the aggregate these individual savings and borrowing should all balance out (and
determine the aggregate interest rate).

Thus, the agents problem is to chose consumption $c$ and assets $a$ to maximize
\begin{align*}
 E[\sum_{t=0}^{\infty} \beta^t \frac{c^{1-\sigma}}{1-\sigma}]
\end{align*}
subject to a budget constraint and no-borrowing constraint
\begin{align*}
 c_t+q a_{t+1}=a_t+e_t \\
 c_t \geq 0, \text{ } a_{t+1} \geq \underline{a}
\end{align*}
The endowments are assumed to follow a two state markov chain $e$. 
The market clearing condition is thus that $\int a d\mu= 0$, where $\mu$ is the equilibrium measure over agent types.

The code implementing the model is following, loosely speaking it recreates the elements of Tables 1 \& 2 of that paper.


\lstinputlisting[language=Matlab]{./Examples/MatlabCodes/Huggett1993.m}

%\end{document}